\chapter{Opis założeń projektu}

\section{Cel projektu}

Celem projektu jest stworzenie aplikacji desktopowej do zarządzania 
kliniką weterynaryjną. System będzie posiadał bezpieczny mechanizm 
logowania i rejestracji z wykorzystaniem bazy danych, a także autoryzację 
użytkowników na podstawie przypisanej do konta roli, takiej jak administrator 
lub użytkownik. Po zalogowaniu aplikacja automatycznie wyświetli odpowiedni 
interfejs, w zależności od roli zalogowanego użytkownika. Administrator 
będzie miał możliwość dodawania i usuwania weterynarzy oraz zabiegów, a 
także przeglądania listy wszystkich zarejestrowanych użytkowników oraz 
złożonych rezerwacji. Użytkownik natomiast będzie mógł doładowywać środki 
na swoje konto, rejestrować swoich pupili, przeglądać historię wizyt oraz 
rezerwować zabiegi u wybranych weterynarzy z uwzględnieniem ich kalendarza 
dostępności.

\vspace{10pt}

\textbf{Przegląd istniejących rozwiązań}

Na rynku dostępnych jest kilka rozwiązań o podobnym przeznaczeniu:

\begin{enumerate}
\item \textbf{VetApp} (\url{https://www.vetapp.app/}) - aplikacja
dedykowana dla gabinetów weterynaryjnych. Oferuje funkcje zarządzania 
wizytami, kartoteką pacjentów oraz fakturowaniem. Jest rozwiązaniem 
płatnym, skierowanym głównie do dużych placówek.
\item \textbf{Vetster} (\url{https://vetster.com/}) - platforma umożliwiająca 
rezerwację wizyt weterynaryjnych online. Jest rozwiązaniem wymagającym opłat subskrypcyjnych.
\item\textbf{eVetPractice} (\url{https://evetpractice.com}) - system 
zarządzania praktyką weterynaryjną oferujący moduły rezerwacji 
i dokumentacji medycznej. Działa w modelu wiążącym się 
z miesięcznymi opłatami i uzależnieniem od infrastruktury dostawcy.
\end{enumerate}

\vspace{5pt}
Przedstawione rozwiązania są zazwyczaj płatne, wymagają dostępu do 
internetu lub są skierowane do dużych placówek. Projektowana aplikacja 
stanowi prostą, bezpłatną alternatywę, dedykowaną 
dla małych i średnich klinik weterynaryjnych.

\section{Realizacja projektu}

Projekt został zrealizowany z wykorzystaniem języka C\# oraz 
technologii WPF (Windows Presentation Foundation). Do przechowywania 
danych użyto relacyjnej bazy danych Microsoft SQL Server 2025, wraz 
z biblioteką ADO.NET i sterownikiem 
Microsoft.Data.SqlClient, która obsługuje informacje o użytkownikach, 
pupilach, weterynarzach, zabiegach oraz rezerwacjach. Aplikacja została 
opracowana w środowisku Visual Studio 2026 Insiders.

\section{Wymagania funkcjonalne oraz niefunkcjonalne}

\textbf{Wymagania funkcjonalne:}
\begin{enumerate}
\item \textbf{Logowanie i rejestracja użytkowników} - aplikacja umożliwia rejestrację nowych użytkowników poprzez formularz, w którym należy podać imię, nazwisko, adres e-mail, numer telefonu, opcjonalny adres zamieszkania oraz hasło. Podczas logowania użytkownik wprowadza swój adres e-mail oraz hasło w celu uwierzytelnienia.

\item \textbf{Panel użytkownika}:
\begin{enumerate}
\item Dodawanie środków do konta - użytkownik może doładować saldo swojego konta, aby móc opłacać rezerwowane zabiegi.
\item Zarządzanie pupilami - użytkownik może dodawać swoich pupili (psy i koty), podając imię, gatunek, rasę, wiek oraz płeć. Może również edytować wiek pupila lub usunąć go z listy.
\item Rezerwacja zabiegu - użytkownik wybiera pupila, następnie zabieg, weterynarza wykonującego dany zabieg, a na końcu wolny termin z kalendarza dostępności wybranego weterynarza. System automatycznie sprawdza dostępność terminów oraz weryfikuje, czy saldo użytkownika jest wystarczające na pokrycie kosztu zabiegu.
\item Przeglądanie zabiegów - użytkownik ma dostęp do historii swoich rezerwacji wraz z informacjami o dacie, godzinie, weterynarzu, zabiegu oraz aktualnym statusie wizyty. Użytkownik może anulować zarezerwowany zabieg.
\item Zarządzanie kontem - użytkownik może edytować swoje dane osobowe oraz zmieniać hasło.
\end{enumerate}

\item \textbf{Panel administratora:}
\begin{enumerate}
\item Zarządzanie użytkownikami - administrator ma dostęp do listy wszystkich zarejestrowanych użytkowników wraz z ich danymi i saldem konta. Może usunąć użytkownika wraz z jego pupilami i rezerwacjami.
\item Zarządzanie rezerwacjami - administrator może przeglądać listę wszystkich złożonych rezerwacji z możliwością filtrowania po dacie, weterynarzu oraz statusie. Administrator może zmieniać status rezerwacji na: Zarezerwowany, Zrealizowany lub Anulowany.
\item Zarządzanie weterynarzami - administrator może dodawać nowych weterynarzy, podając imię, nazwisko, specjalizację, email, telefon, hasło oraz godziny pracy w poszczególnych dniach tygodnia. Może również przypisywać zabiegi do weterynarza oraz usuwać weterynarzy z systemu.
\item Zarządzanie zabiegami - administrator może dodawać nowe zabiegi do oferty kliniki, wprowadzając nazwę, czas trwania oraz cenę. Może także usuwać dostępne zabiegi.
\end{enumerate}
\item Opcja wylogowania - zarówno panel użytkownika, jak i panel administratora, posiadają opcję wylogowania, która przenosi z powrotem do ekranu logowania.
\end{enumerate}
\vspace{5pt}

\textbf{Wymagania niefunkcjonalne:}
\begin{enumerate}
\item Bezpieczeństwo - dostęp do poszczególnych funkcji ograniczony zgodnie z rolą: użytkownik, administrator. Walidacja danych wejściowych na poziomie aplikacji zapobiega wprowadzaniu nieprawidłowych danych.
\item Skalowalność - struktura bazy danych oraz logika aplikacji umożliwiają rozszerzenie systemu o dodatkowe funkcjonalności, np. obsługę kolejnych gatunków zwierząt, powiadomienia o wizytach lub wystawianie faktur.
\item Responsywność - czas odpowiedzi programu na operacje użytkownika nie powinien przekraczać 1 sekundy.
\item Intuicyjność - program łatwy do użytku dzięki przejrzystej oprawie graficznej oraz podziałowi na logiczne sekcje nawigacyjne.
\end{enumerate}